% Source: https://artofproblemsolving.com/community/c3h1778586_math_problems_week
% Stable identity: aops-topic-1778586-q02
A square inscribed in a circle of radius $r$ has diagonal $2r$, so its area is $\dfrac{(2r)^2}{2}=2r^2$, and its inscribed circle has radius equal to half the side, $\dfrac{r}{\sqrt2}$.

\begin{asy}
usepackage("amsmath");
size(5cm);
defaultpen(fontsize(10pt));
real r=1;
for(int k=0;k<5;++k){
  fill(circle((0,0),r),lightgray);
  path sq=r*dir(45)--r*dir(135)--r*dir(225)--r*dir(315)--cycle;
  fill(sq,white);
  draw(circle((0,0),r),linewidth(0.7)); draw(sq,linewidth(0.7));
  r=r/sqrt(2);
}
draw((0,0)--dir(160),blue+linewidth(0.6)); label("$1$",0.88*dir(160)+0.13*dir(70),blue);
label("\mbox{}",(min(currentpicture,user=true)+max(currentpicture,user=true))/2,black); // keeps the text after the figure black when LaTeX typesets the labels inline
\end{asy}

So the shaded region at a stage with circle radius $r$ has area $\pi r^2-2r^2=(\pi-2)r^2$, and the next circle has radius $\dfrac {r}{\sqrt2}$, i.e.\ $r^2$ is halved at each stage. Starting from $r=1$, the total is the geometric series
\[
(\pi-2)\left(1+\dfrac{1}{2}+\dfrac{1}{4}+\cdots\right)=(\pi-2)\cdot2=\boxed{2\pi-4}.
\]
